\Needspace{35\baselineskip}\SourceLeaf{84}

\begin{center}
I.
\end{center}

\noindent\hspace*{\parindent}Da für $u = 0$ die Gleichung $x = t$ gilt, folgt nach dem Lehrsatz von \textsc{Taylor}

\[
\varphi x = \varphi t + u \frac{d\varphi x}{du} + \frac{uu}{1.2}\frac{d^2\varphi x}{du^2}
+ \frac{u^3}{1.2.3}\frac{d^3\varphi x}{du^3} + \text{etc.},
\]

\Needspace{23\baselineskip}\noindent wobei die Werte der Größen $\dfrac{d\varphi x}{du}$, $\dfrac{d^2\varphi x}{du^2}$ usw.\ für
$u = 0$ einzusetzen sind. Allgemein gilt jedoch

\[
\frac{d\varphi x}{du} = X\frac{d\varphi x}{dt},
\]

\[
\frac{d^2\varphi x}{du^2} = \frac{dX^2 \frac{d\varphi x}{dt}}{dt},
\]

\[
\frac{d^3\varphi x}{du^3} = \frac{d^2X^3 \frac{d\varphi x}{dt}}{dt^2},
\]

\[
\cdot\quad\cdot\quad\cdot\quad\cdot\quad\cdot\quad\cdot\quad\cdot\quad\cdot\quad\cdot,
\]

\noindent wie ich sogleich beweisen werde. Ferner sieht man leicht, dass für $u = 0$, d.\,h.\
$x = t$, diese Ausdrücke

\[
X\frac{d\varphi x}{dt},\quad \frac{dX^2 \frac{d\varphi x}{dt}}{dt},\quad
\text{etc.\ in}\quad X\frac{d\varphi x}{dx},\quad \frac{dX^2 \frac{d\varphi x}{dx}}{dx},\quad \text{etc.}
\]

\noindent übergehen, wobei dann $x = t$ einzusetzen ist. Daraus ergibt sich also

\[
\varphi x = \varphi t + uX\frac{d\varphi x}{dx} + \frac{uu}{1.2}\frac{dX^2 \frac{d\varphi x}{dx}}{dx}
+ \frac{u^3}{1.2.3}\frac{d^2X^3 \frac{d\varphi x}{dx}}{dx^2} + \cdot\,\cdot\,\cdot,
\]

\noindent wobei $x = t$ einzusetzen ist. Das ist die Formel des Herrn \textsc{de la Grange}.

\medskip

\Needspace{10\baselineskip}\begin{center}
II.
\end{center}

\noindent\hspace*{\parindent}Die hierbei verwendeten Umformungen beweise ich wie folgt:

\noindent\hspace*{\parindent}Ich differenziere die gegebene Gleichung $x = t + uX$ nach $t$ und nach $u$
und erhalte daraus

\[
\frac{dx}{dt} = 1 + u\frac{dX}{dt} = 1 + u\frac{dX}{dx}\frac{dx}{dt},
\]

\[
\frac{dx}{du} = X + u\frac{dX}{du} = X + u\frac{dX}{dx}\frac{dx}{du}.
\]
